% ============================================================
%  PREAMBOLO COMUNE — Pacchetti e Colori
%  Da includere nel documento principale prima di \begin{document}
% ============================================================

% --- Pacchetti ---
\usepackage{tikz}
\usepackage{amsmath}
\usetikzlibrary{
  arrows.meta,
  calc,
  positioning,
  decorations.markings,
  decorations.pathmorphing,
  patterns
}

% ============================================================
%  PALETTE COLORI UNIFICATA  (un solo nome per ciascun colore)
% ============================================================

% ---- Rossi / Rosa ----
\definecolor{posred}{RGB}{239,130,140}      % rosa‐rosso per cariche positive (+)
\definecolor{ered}{RGB}{227,35,35}          % rosso vivo per vettori E, F, v, frecce campo

% ---- Blu / Azzurri ----
\definecolor{negblue}{RGB}{120,190,230}     % azzurro per cariche negative (-)
\definecolor{surfblue}{RGB}{205,230,245}    % azzurro chiaro per superfici chiuse, Gauss
\definecolor{bblue}{RGB}{85,155,220}        % blu per vettori B, H, D, linee di flusso magnetico

% ---- Viola ----
\definecolor{vpurple}{RGB}{160,80,160}      % viola per vettori p, M, m, j, correnti

% ---- Arancioni / Beige ----
\definecolor{charge}{RGB}{245,175,120}      % arancione per sfere di carica puntiforme
\definecolor{blobpink}{RGB}{245,210,195}    % rosa pallido per corpi estesi / distribuzioni
\definecolor{ytangent}{RGB}{250,235,190}    % giallo paglierino per aree tangenti / parallelogrammi

% ---- Grigi ----
\definecolor{blobgray}{RGB}{220,220,220}    % grigio chiaro per conduttori / corpi metallici
\definecolor{planegray}{RGB}{230,230,230}   % grigio chiarissimo per piani di separazione

% ============================================================
%  COLORI AGGIUNTIVI (usati solo nei file preesistenti
%  image1–4, image55, image63 e da NON rinominare)
% ============================================================
%  I file image55.tex e image63.tex, già scritti e non modificabili,
%  usano i seguenti alias che sono identici a colori della palette:
%
%    matpink   = blobpink      RGB(245,210,195)
%    matpinkdk = blobpink scuro  RGB(225,175,155)   (solo image55)
%    loopblue  = bblue variant   RGB(90,130,200)    (solo image55)
%    mred      = ered            RGB(227,35,35)     (solo image55)
%    surfedge  = bblue variant   RGB(80,120,170)    (solo image63)
%    ghost     = grigio medio    RGB(140,140,140)   (solo image63)
%    vred      = ered            RGB(227,35,35)     (solo image63)
%    cubecol   = blobpink tenue  RGB(235,190,170)   (solo image3)
%
%  Se si compilano quei file all'interno dello stesso documento,
%  aggiungere queste definizioni:
\definecolor{matpink}{RGB}{245,210,195}
\definecolor{matpinkdk}{RGB}{225,175,155}
\definecolor{loopblue}{RGB}{90,130,200}
\definecolor{mred}{RGB}{227,35,35}
\definecolor{surfedge}{RGB}{80,120,170}
\definecolor{ghost}{RGB}{140,140,140}
\definecolor{vred}{RGB}{227,35,35}
\definecolor{cubecol}{RGB}{235,190,170}
